% --- Archon physics formalization source begin ---
% archon:physics
% archon:covers PhyXMiniProblems/problem_phyx_mini_0105.lean
% archon:source-report reports/phyx_mini/problem_phyx_mini_0105.source.json

\chapter{Physics problem phyx\_mini\_0105}
\label{ch:PhyXMiniProblems_problem_phyx_mini_0105}

\paragraph{Problem source.}
\#\# Physical scenario

A beam of light traveling horizontally is made of an unpolarized component with intensity \$I\_0\$ and a polarized component with intensity \$I\_p\$. The plane of polarization of the polarized component is oriented at an angle \$\textbackslash{}theta\$ with respect to the vertical. Figure is a graph of the total intensity \$I\_\{\textbackslash{}text\{total\}\}\$ after the light passes through a polarizer versus the angle \$\textbackslash{}alpha\$ that the polarizer's axis makes with respect to the vertical.

\#\# Question

What are the values of \$I\_p\$?

\#\# Answer choices

- A: A: 20W/m\textasciicircum{}2

- B: B: 15W/m\textasciicircum{}2

- C: C: 25W/m\textasciicircum{}2

- D: D: 30W/m\textasciicircum{}2

\#\# Figure caption (auxiliary; use the image as primary evidence)

The image is a scatter plot graph.

- The x-axis is labeled as "α (°)" and ranges from 0 to 200.
- The y-axis is labeled as "Iₜₒₜₐₗ (W/m²)" and ranges from 0 to 30.
- The data points form a sinusoidal pattern, with peaks around 25 W/m² and troughs around 10 W/m².
- The oscillation appears approximately periodic with peaks and troughs at around 50° and 150°, respectively.
- There are grid lines for better visibility of data points.

The graph likely illustrates a relationship between angle (α) and intensity (Iₜₒₜₐₗ).

\#\# Dataset metadata

Wave Optics; Physical Model Grounding Reasoning, Multi-Formula Reasoning

\paragraph{Recorded answer/context.}
A: A: 20W/m\textasciicircum{}2

\paragraph{Figure/image path.}
/root/proposal\_for\_physic/science-mango/phyx\_mini\_run/phyx\_data/test\_image/105.png

\paragraph{Formalization target.}
create a compiling Lean file with sorry bodies at `PhyXMiniProblems/problem\_phyx\_mini\_0105.lean`. The Lean declarations must preserve the physical quantities, dimensions or dimensional roles, figure labels, governing-law hypotheses, and final relation expressed by this problem.
Use Mathlib/Physlib names found through LeanExplore where available. If a physics API is missing, introduce faithful local abstractions rather than scalar placeholder aliases.

\begin{theorem}[Physics formalization target]
\label{thm:physics:phyx_mini_0105:target}
\lean{PhyXMiniProblems.ProblemPhyXMini0105.problem_phyx_mini_0105}
\uses{def:physics:phyx-mini-0105:phyxminiproblems-problemphyxmini0105-wattspersquaremeter, def:physics:phyx-mini-0105:phyxminiproblems-problemphyxmini0105-mixedbeampolarizersetup, def:physics:phyx-mini-0105:phyxminiproblems-problemphyxmini0105-matchesproblemdescription, def:physics:phyx-mini-0105:phyxminiproblems-problemphyxmini0105-satisfiesidealpolarizerlaws, def:physics:phyx-mini-0105:phyxminiproblems-problemphyxmini0105-matchesintensitygraph, def:physics:phyx-mini-0105:phyxminiproblems-problemphyxmini0105-matchespolarizedintensityanswer}
The assigned autoformalize agent should translate this physics problem into Lean declarations in the covered file, with theorem and lemma proof bodies written as `by sorry`.
\end{theorem}
\begin{proof}
This is an autoformalization task, not a proof task. Produce statements that can later be proved by the physics prover without weakening the physical model.
\end{proof}

% --- Archon declaration topology begin ---
\section{Declaration topology}
\begin{definition}[Irradiance Quantity]
  \label{def:physics:phyx-mini-0105:phyxminiproblems-problemphyxmini0105-irradiancequantity}
  \lean{PhyXMiniProblems.ProblemPhyXMini0105.IrradianceQuantity}
  A nonnegative physical optical irradiance (power per area)
\end{definition}
\begin{proof}
  This is the declared model definition of Irradiance Quantity.
\end{proof}
\begin{definition}[irradiance In Watts Per Square Meter]
  \label{def:physics:phyx-mini-0105:phyxminiproblems-problemphyxmini0105-irradianceinwattspersquaremeter}
  \lean{PhyXMiniProblems.ProblemPhyXMini0105.irradianceInWattsPerSquareMeter}
  \uses{def:physics:phyx-mini-0105:phyxminiproblems-problemphyxmini0105-irradiancequantity}
  The scalar SI readout of an irradiance, in watts per square meter
\end{definition}
\begin{proof}
  This is the declared model definition of irradiance In Watts Per Square Meter.
\end{proof}
\begin{definition}[watts Per Square Meter]
  \label{def:physics:phyx-mini-0105:phyxminiproblems-problemphyxmini0105-wattspersquaremeter}
  \lean{PhyXMiniProblems.ProblemPhyXMini0105.wattsPerSquareMeter}
  \uses{def:physics:phyx-mini-0105:phyxminiproblems-problemphyxmini0105-irradiancequantity}
  The dimensionful irradiance whose SI readout is the supplied number
\end{definition}
\begin{proof}
  This is the declared model definition of watts Per Square Meter.
\end{proof}
\begin{definition}[degrees]
  \label{def:physics:phyx-mini-0105:phyxminiproblems-problemphyxmini0105-degrees}
  \lean{PhyXMiniProblems.ProblemPhyXMini0105.degrees}
  Convert a numerical degree coordinate from the graph into a physical angle
\end{definition}
\begin{proof}
  This is the declared model definition of degrees.
\end{proof}
\begin{definition}[Propagation Direction]
  \label{def:physics:phyx-mini-0105:phyxminiproblems-problemphyxmini0105-propagationdirection}
  \lean{PhyXMiniProblems.ProblemPhyXMini0105.PropagationDirection}
  Propagation direction stated in the problem text
\end{definition}
\begin{proof}
  This is the declared model definition of Propagation Direction.
\end{proof}
\begin{definition}[Figure Axis]
  \label{def:physics:phyx-mini-0105:phyxminiproblems-problemphyxmini0105-figureaxis}
  \lean{PhyXMiniProblems.ProblemPhyXMini0105.FigureAxis}
  The two labeled axes in the supplied scatter plot
\end{definition}
\begin{proof}
  This is the declared model definition of Figure Axis.
\end{proof}
\begin{definition}[Figure Axis Unit]
  \label{def:physics:phyx-mini-0105:phyxminiproblems-problemphyxmini0105-figureaxisunit}
  \lean{PhyXMiniProblems.ProblemPhyXMini0105.FigureAxisUnit}
  \uses{def:physics:phyx-mini-0105:phyxminiproblems-problemphyxmini0105-wattspersquaremeter, def:physics:phyx-mini-0105:phyxminiproblems-problemphyxmini0105-degrees}
  The units printed beside the two graph axes
\end{definition}
\begin{proof}
  This is the declared model definition of Figure Axis Unit.
\end{proof}
\begin{definition}[Polarizer Intensity Graph]
  \label{def:physics:phyx-mini-0105:phyxminiproblems-problemphyxmini0105-polarizerintensitygraph}
  \lean{PhyXMiniProblems.ProblemPhyXMini0105.PolarizerIntensityGraph}
  \uses{def:physics:phyx-mini-0105:phyxminiproblems-problemphyxmini0105-figureaxis, def:physics:phyx-mini-0105:phyxminiproblems-problemphyxmini0105-figureaxisunit}
  Metadata and plotted ranges visible in the source graph
\end{definition}
\begin{proof}
  This is the declared model definition of Polarizer Intensity Graph.
\end{proof}
\begin{definition}[Mixed Beam Polarizer Setup]
  \label{def:physics:phyx-mini-0105:phyxminiproblems-problemphyxmini0105-mixedbeampolarizersetup}
  \lean{PhyXMiniProblems.ProblemPhyXMini0105.MixedBeamPolarizerSetup}
  \uses{def:physics:phyx-mini-0105:phyxminiproblems-problemphyxmini0105-irradiancequantity, def:physics:phyx-mini-0105:phyxminiproblems-problemphyxmini0105-propagationdirection, def:physics:phyx-mini-0105:phyxminiproblems-problemphyxmini0105-polarizerintensitygraph}
  The physical components and figure-labeled quantities. The arguments of the three transmitted-irradiance functions are the plotted α coordinates in degrees. The polarized and unpolarized transmitted components are retained separately so that the two governing transmission laws are explicit
\end{definition}
\begin{proof}
  This is the declared model definition of Mixed Beam Polarizer Setup.
\end{proof}
\begin{definition}[Matches Problem Description]
  \label{def:physics:phyx-mini-0105:phyxminiproblems-problemphyxmini0105-matchesproblemdescription}
  \lean{PhyXMiniProblems.ProblemPhyXMini0105.MatchesProblemDescription}
  \uses{def:physics:phyx-mini-0105:phyxminiproblems-problemphyxmini0105-mixedbeampolarizersetup}
  The textual setup fact that the incident beam travels horizontally
\end{definition}
\begin{proof}
  This is the declared model definition of Matches Problem Description.
\end{proof}
\begin{definition}[Satisfies Ideal Polarizer Laws]
  \label{def:physics:phyx-mini-0105:phyxminiproblems-problemphyxmini0105-satisfiesidealpolarizerlaws}
  \lean{PhyXMiniProblems.ProblemPhyXMini0105.SatisfiesIdealPolarizerLaws}
  \uses{def:physics:phyx-mini-0105:phyxminiproblems-problemphyxmini0105-irradianceinwattspersquaremeter, def:physics:phyx-mini-0105:phyxminiproblems-problemphyxmini0105-degrees, def:physics:phyx-mini-0105:phyxminiproblems-problemphyxmini0105-mixedbeampolarizersetup}
  The ideal-polarizer laws used by the problem. An ideal polarizer transmits half of an unpolarized component, a linearly polarized component obeys Malus's law, and the two mutually incoherent transmitted irradiances add. These laws are stated for every polarizer angle and contain no numerical value for Iₚ
\end{definition}
\begin{proof}
  This is the declared model definition of Satisfies Ideal Polarizer Laws.
\end{proof}
\begin{definition}[Matches Intensity Graph]
  \label{def:physics:phyx-mini-0105:phyxminiproblems-problemphyxmini0105-matchesintensitygraph}
  \lean{PhyXMiniProblems.ProblemPhyXMini0105.MatchesIntensityGraph}
  \uses{def:physics:phyx-mini-0105:phyxminiproblems-problemphyxmini0105-irradianceinwattspersquaremeter, def:physics:phyx-mini-0105:phyxminiproblems-problemphyxmini0105-wattspersquaremeter, def:physics:phyx-mini-0105:phyxminiproblems-problemphyxmini0105-degrees, def:physics:phyx-mini-0105:phyxminiproblems-problemphyxmini0105-mixedbeampolarizersetup}
  Calibrated readouts from the primary figure. Its horizontal axis is α in degrees from 0 to 200; its vertical axis is total irradiance in W/m² from 0 to 30. On that plotted interval the trace attains a maximum of 25 W/m² and a minimum of 5 W/m². The auxiliary generated caption's claim that the trough is near 10 W/m² conflicts with the plotted points; the primary image is used as required. Neither extremum assumption states the requested polarized-component irradiance
\end{definition}
\begin{proof}
  This is the declared model definition of Matches Intensity Graph.
\end{proof}
\begin{definition}[Answer Choice]
  \label{def:physics:phyx-mini-0105:phyxminiproblems-problemphyxmini0105-answerchoice}
  \lean{PhyXMiniProblems.ProblemPhyXMini0105.AnswerChoice}
  Labels of the four answer choices printed with the problem
\end{definition}
\begin{proof}
  This is the declared model definition of Answer Choice.
\end{proof}
\begin{definition}[answer Irradiance In Watts Per Square Meter]
  \label{def:physics:phyx-mini-0105:phyxminiproblems-problemphyxmini0105-answerirradianceinwattspersquaremeter}
  \lean{PhyXMiniProblems.ProblemPhyXMini0105.answerIrradianceInWattsPerSquareMeter}
  \uses{def:physics:phyx-mini-0105:phyxminiproblems-problemphyxmini0105-answerchoice}
  The irradiance readout in W/m² printed beside each answer choice
\end{definition}
\begin{proof}
  This is the declared model definition of answer Irradiance In Watts Per Square Meter.
\end{proof}
\begin{definition}[Matches Polarized Intensity Answer]
  \label{def:physics:phyx-mini-0105:phyxminiproblems-problemphyxmini0105-matchespolarizedintensityanswer}
  \lean{PhyXMiniProblems.ProblemPhyXMini0105.MatchesPolarizedIntensityAnswer}
  \uses{def:physics:phyx-mini-0105:phyxminiproblems-problemphyxmini0105-irradianceinwattspersquaremeter, def:physics:phyx-mini-0105:phyxminiproblems-problemphyxmini0105-mixedbeampolarizersetup, def:physics:phyx-mini-0105:phyxminiproblems-problemphyxmini0105-answerchoice, def:physics:phyx-mini-0105:phyxminiproblems-problemphyxmini0105-answerirradianceinwattspersquaremeter}
  A choice matches the SI readout of the incident polarized component
\end{definition}
\begin{proof}
  This is the declared model definition of Matches Polarized Intensity Answer.
\end{proof}
% --- Archon declaration topology end ---
% --- Archon physics formalization source end ---
